% =====================================================================
%  2bat-ud6-9.tex  —  SESSIÓ 9: Significació estadística: consolidació
%  (sense preàmbul: s'inclou des de main.tex)
% =====================================================================
\horatitol{Sessió 9 — Significació estadística: consolidació}%
{Assentar la idea de significació amb una caça de l'error sobre conclusions precipitades i una bateria de casos, amb full d'autoavaluació.}

\subsection{Idea clau}

Consolidem una de les idees més valuoses de la unitat: la \textbf{significació
estadística}. Ho fem amb una \textbf{caça de l'error} sobre conclusions precipitades i una
bateria de casos on cal valorar si una diferència és real o casual, tenint en compte la
mida de la diferència, la mida de les mostres i la superposició dels intervals.

\begin{idea}
\textbf{Una diferència és significativa si és difícilment atribuïble a l'atzar.}
\begin{itemize}[leftmargin=1.4em, itemsep=2pt]
  \item \textbf{Conclusió precipitada:} afirmar que un grup és millor que un altre a partir
        d'una diferència \emph{petita} amb mostres \emph{petites}.
  \item \textbf{Conclusió prudent:} afirmar-ho només si la diferència és prou gran i les
        mostres prou grans perquè els intervals \textbf{no} se superposin.
\end{itemize}
\textbf{Regla pràctica:} si els intervals de confiança dels dos grups \textbf{se
superposen}, la diferència no es pot considerar significativa.
\end{idea}

\subsection{Exemples}

\begin{exemple}[caça l'error: una conclusió precipitada]
Afirmació: «El programa A ($52\%$) és millor que el B ($50\%$)», a partir de mostres de $40$
persones cadascuna. \textbf{Error:} $2$ punts de diferència amb mostres tan petites entren
de ple dins del marge d'error; els intervals se superposen àmpliament. La conclusió
\textbf{no} està justificada: la diferència és perfectament atribuïble a l'atzar.
\end{exemple}

\begin{exemple}[una conclusió ben feta]
Afirmació: «El programa A ($64\%$) supera el B ($48\%$)», amb mostres de $800$ persones i
marges de $\pm 3\%$. Intervals $[61\%,67\%]$ i $[45\%,51\%]$: \textbf{no} se superposen, i
la diferència és de $16$ punts. Aquí la conclusió \textbf{sí} està justificada: la
diferència és significativa.
\end{exemple}

\subsection{Exercicis resolts}

\begin{exresolt}[detectar la conclusió precipitada]
«L'edició d'enguany ($57\%$) ha estat millor que la de l'any passat ($55\%$)», amb mostres
de $50$ persones. Està justificada la conclusió?
\solucio
\textbf{No.} Amb mostres de $50$ persones el marge d'error és gran, així que $2$ punts de
diferència queden dins del que pot explicar l'atzar; els intervals se superposarien. La
conclusió és \textbf{precipitada}: no es pot afirmar que hagi estat millor.
\resposta{No està justificada: diferència petita i mostres petites.}
\end{exresolt}

\begin{exresolt}[valorar amb tots els factors]
Dos grups donen un $40\%\pm 3\%$ i un $50\%\pm 3\%$. Valora si la diferència és
significativa considerant la diferència i els intervals.
\solucio
Intervals: $[37\%,43\%]$ i $[47\%,53\%]$. \textbf{No} se superposen (hi ha un buit del
$43\%$ al $47\%$) i la diferència és de $10$ punts. La diferència sembla
\textbf{significativa}: difícilment s'explica només per l'atzar.
\resposta{Significativa: intervals separats i diferència de $10$ punts.}
\end{exresolt}

\subsection{Exercicis per fer}

\noindent\textit{Itinerari mínim: casos amb diferències clarament grans o clarament
negligibles (ex. 1--3). Ampliació: casos ambigus (ex. 4--6).}

\begin{exercicis}
\begin{exlist}
  \item \textbf{(Mínim)} Dos grups: $30\%\pm 3\%$ i $55\%\pm 3\%$. Se superposen els
        intervals? La diferència és significativa?
  \item \textbf{(Mínim)} Dos grups: $49\%\pm 4\%$ i $51\%\pm 4\%$. Se superposen? La
        diferència és significativa?
  \item \textbf{(Mínim, caça l'error)} «A ($60\%$) és molt millor que B ($59\%$)», amb
        mostres de $35$ persones. Justifica per què la conclusió és precipitada.
  \item \textbf{(Ampliació)} Dos grups: $52\%\pm 4\%$ i $59\%\pm 4\%$. Se superposen els
        intervals? Es pot concloure alguna cosa amb seguretat?
  \item \textbf{(Ampliació)} La mateixa diferència de $7$ punts, amb marges de $\pm 2\%$ en
        lloc de $\pm 4\%$. Ara se superposen els intervals? Canvia la conclusió?
  \item \textbf{(Ampliació, raonament)} Explica com pot ser que una diferència de $5$ punts
        sigui significativa en un estudi (mostres grans) i no ho sigui en un altre (mostres
        petites).
\end{exlist}
\end{exercicis}

% ---------------------------------------------------------------------
\begin{idea}
\textbf{Full d'autoavaluació.} Marca amb sinceritat com et trobes, de cara a la prova:
\begin{center}
\renewcommand{\arraystretch}{1.4}
\begin{tabular}{p{8.4cm} c c c}
\toprule
\textbf{Sé fer\dots} & \textbf{Sol} & \textbf{Amb ajuda} & \textbf{Encara no} \\
\midrule
Construir l'interval de confiança d'un resultat & & & \\
Veure si dos intervals se superposen & & & \\
Valorar si una diferència és significativa o casual & & & \\
Detectar una conclusió precipitada & & & \\
Tenir en compte la mida de les mostres & & & \\
\bottomrule
\end{tabular}
\end{center}
\end{idea}

% ---------------------------------------------------------------------
\fullsolucions{Sessió 9}
\begin{solucions}
\begin{sollist}
  \item Intervals $[27\%,33\%]$ i $[52\%,58\%]$: \textbf{no} se superposen. La diferència
        de $25$ punts és \textbf{significativa}.
  \item Intervals $[45\%,53\%]$ i $[47\%,55\%]$: \textbf{se superposen} àmpliament. La
        diferència de $2$ punts \textbf{no} és significativa.
  \item Amb mostres de $35$ i $1$ punt de diferència, el marge d'error és gran i els
        intervals se superposen del tot; la diferència és perfectament atzar. La conclusió
        «molt millor» és precipitada.
  \item Intervals $[48\%,56\%]$ i $[55\%,63\%]$: \textbf{se superposen} lleugerament (del
        $55\%$ al $56\%$). No es pot concloure amb seguretat que un sigui millor; és un cas
        ambigu.
  \item Amb $\pm 2\%$: intervals $[50\%,54\%]$ i $[57\%,61\%]$: \textbf{no} se superposen.
        Ara sí que la diferència de $7$ punts es pot considerar significativa. Reduir el
        marge (mostres més grans) canvia la conclusió.
  \item Resposta oberta. Idea: amb mostres grans el marge és petit i els intervals estrets,
        de manera que $5$ punts poden separar-los (significativa); amb mostres petites el
        marge és gran i els intervals amplis se superposen, de manera que els mateixos $5$
        punts poden ser atzar (no significativa).
\end{sollist}
\end{solucions}
